
\documentclass{article}
\usepackage{colortbl}
\usepackage{makecell}
\usepackage{multirow}
\usepackage{supertabular}

\begin{document}

\newcounter{utterance}

\centering \large Interaction Transcript for game `textmapworld', experiment `small', episode 0 with reka{-}flash{-}3.1{-}t0.0.
\vspace{24pt}

{ \footnotesize  \setcounter{utterance}{1}
\setlength{\tabcolsep}{0pt}
\begin{supertabular}{c@{$\;$}|p{.15\linewidth}@{}p{.15\linewidth}p{.15\linewidth}p{.15\linewidth}p{.15\linewidth}p{.15\linewidth}}
    \# & $\;$A & \multicolumn{4}{c}{Game Master} & $\;\:$B\\
    \hline

    \theutterance \stepcounter{utterance}  
    & & \multicolumn{4}{p{0.6\linewidth}}{
        \cellcolor[rgb]{0.9,0.9,0.9}{
            \makecell[{{p{\linewidth}}}]{
                \texttt{\tiny{[P1$\langle$GM]}}
                \texttt{Please help me with the following task. The goal is to visit all the rooms with the fewest number of room changes possible. In each room, you need to decide the direction to go in. Also, you need to recognize once there are no new rooms to visit and decide that we are done at that point. Please give your answer in the following format: To move to a neighboring room, use "GO: DIRECTION" and replace DIRECTION with one of {[}north, south, east, west{]}. To stop the exploration, answer with "DONE" instead. Omit any other text.} \\
\texttt{Here is an example:} \\
\texttt{You are in the Kitchen. Currently available directions: south, west. What is your next instruction?} \\
\texttt{GO: west} \\
\texttt{You have made a step and entered a Lobby. Currently available directions: east, north. What is your next instruction?} \\
\texttt{GO: north} \\
\texttt{...} \\
\texttt{You have made a step and entered a Bedroom. Currently available directions: south. What is your next instruction?} \\
\texttt{DONE} \\
\texttt{Let us start. You are in the Studio. Currently available directions: west. What is your next instruction?} \\
            }
        }
    }
    & \\ \\

    \theutterance \stepcounter{utterance}  
    & \multicolumn{4}{p{0.6\linewidth}}{
        \cellcolor[rgb]{0.9,0.9,0.9}{
            \makecell[{{p{\linewidth}}}]{
                \texttt{\tiny{[P1$\rangle$GM]}}
                \texttt{<reasoning>} \\
\texttt{Let's tackle this problem. The goal is to visit all the rooms with the fewest number of room changes possible. So, I need to figure out the optimal path that covers all rooms in the least steps. But how do I know the layout of the rooms? The user hasn't provided a map, so I must be working with the information given step by step as I move through the rooms.} \\
\\ 
\texttt{First, the starting point is the Studio, and the available direction is west. So the first move is probably to go west. But wait, maybe there's a strategy here. Since we want to minimize room changes, which I think refers to the number of moves (each time you go from one room to another counts as a change), we need to explore in a way that covers as many rooms as possible with each move. But without knowing the layout, it's tricky. } \\
\\ 
\texttt{Wait, maybe this is similar to a graph traversal problem, like BFS (Breadth{-}First Search) or DFS (Depth{-}First Search). BFS is good for finding the shortest path in an unweighted graph, which might be what's needed here. But since we don't have the entire map, we have to build it as we go. So each time we enter a room, we note the available directions and the rooms connected to them. Then, we can plan a path that covers all rooms with the least steps.} \\
\\ 
\texttt{But how to do that without knowing the entire structure? Maybe we need to keep track of visited rooms and the paths taken. Let's think step by step.} \\
\\ 
\texttt{Starting in the Studio. Available direction: west. Let's go west first. Suppose that takes us to another room, say Room A. From Room A, available directions might be east (back to Studio) and maybe others. But we need to explore all possible directions. However, to minimize room changes, we might want to explore as much as possible in one direction before backtracking. That's similar to DFS, which goes deep before wide. Alternatively, BFS would explore all possible directions level by level. But since we don't have the map, maybe a greedy approach is needed.} \\
\\ 
\texttt{Wait, the problem says "visit all the rooms with the fewest number of room changes possible." So the minimal number of moves (each move is a room change) to visit all rooms. So this is equivalent to finding the shortest path that visits every room at least once. That's the Traveling} \\
            }
        }
    }
    & & \\ \\

    \theutterance \stepcounter{utterance}  
    & & & \multicolumn{2}{p{0.3\linewidth}}{
        \cellcolor[rgb]{0.9,0.9,0.9}{
            \makecell[{{p{\linewidth}}}]{
                \texttt{\tiny{[GM$|$GM]}}
                \texttt{abort game} \\
            }
        }
    }
    & & \\ \\

\end{supertabular}
}

\end{document}
